% Part: first-order-logic
% Chapter: axiomatic-deduction
% Section: axioms-rules-quantifiers

\documentclass[../../../include/open-logic-section]{subfiles}

\begin{document}

\olfileid{fol}{axd}{qua}

\olsection{بديهيات المكمّمات وقواعدها}

\begin{defn}[بديهيات المكمّمات]
\emph{بديهيات} المكمّمات هي كل ما يقع تحت الصورتين الآتيتين:
\begin{align}
\ollabel{ax:q1} & \lforall[x][!B] \lif !B(t), \\
\ollabel{ax:q2} & !B(t) \lif \lexists[x][!B],
\end{align}
حيث يؤخذ~$t$ أيّ حد مغلق.
\end{defn}

% تصحيح مصدر (0118-I1): أُضيف مخزون محافظ غير متناهٍ للثوابت المميزة.
ونلحق بلغة المسألة~$\Lang L$ طائفةً لا تتناهى
$\mathcal E=\{e_0,e_1,\dots\}$ من الثوابت المساعدة المستجدة. وقد تقع ثوابت
$\mathcal E$ في أثناء البرهان، وتُحمل صور البديهيات والقواعد على اللغة
الموسعة، ولا يقع شيء منها في دعم ال!!{derivation} ولا في خاتمته. ومنها
نأخذ الثابت المميّز في كل تطبيق من~$\mathcal E$.

\begin{defn}[قواعد المكمّمات]
% تصحيح مصدر (0115-I1): أُضيف شرط عدم ورود الثابت المميز في A(x)، وهو ساقط من الأصل الإنجليزي وخط الأساس العربي.
\begin{enumerate}
\item متى سبقت~$!B \lif !A(a)$ في ال!!{derivation}، وكان
  $a$ غير وارد في دعم ال!!{derivation} ولا في~$!B$ ولا في~$!A(x)$، جاز إلحاق~$!B \lif
  \lforall[x][!A(x)]$ خطوةَ استدلال صحيحة.
\item ومتى سبقت~$!A(a) \lif !B$ في ال!!{derivation}، وكان
  $a$ غير وارد في دعم ال!!{derivation} ولا في~$!B$ ولا في~$!A(x)$، جاز إلحاق~$\lexists[x][!A(x)] \lif
  !B$ خطوةَ استدلال صحيحة.
\end{enumerate}
\end{defn}

ونرمز إلى أيّ من هاتين القاعدتين اختصارًا بـ«\QR».

ونعد من الثوابت الواقعة في شجرة برهان كل ثابت يقع في صيغة منها، وكل ثابت
مميز مكتوب على عقدة~\QR{}، وإن كان المتغير المسوّر غير واقع في الصيغة فلم
يظهر ذلك الثابت في مقدمة العقدة.

\begin{lem}[تجديد الثوابت المميزة]\ollabel{lem:qr-freshening}
لتكن~$\Pi$ !!a{derivation} متناهية الدعم خاتمتها~$!C$، ولتكن~$F$
طائفة متناهية من الثوابت. توجد !!a{derivation}~$\Pi'$ خاتمتها~$!C$،
ودعمها داخل في دعم~$\Pi$، وتكون الثوابت المميزة في تطبيقات~\QR{} فيها
متمايزة، ولا يقع شيء منها في~$F$ ولا في دعم~$\Pi$ ولا في~$!B$ أو
~$!A(x)$ الخاصين بتطبيقها.
\end{lem}

\begin{proof}
سم دعم~$\Pi$ الأول~$\Sigma_0$، ثم احذف سطوره التي لا تستعمل، وابسط
الباقي إلى شجرة برهان متناهية دعمها~$\Sigma\subseteq\Sigma_0$.
ونستقرئ في العبارة الأقوى الآتية: إذا كانت~$T$ شجرة متناهية جذرها~$!D$
ودعمها~$\Sigma$، وكانت~$H$ طائفة متناهية من الثوابت، أمكن إنشاء
شجرة~$T'$ لها الجذر والدعم نفسهما، وثوابتها المميزة متباينة، لا يقع شيء منها في~$H$ ولا في
$!D$ ولا في~$\Sigma$، مع بقاء جميع شروط~\QR{} المحلية. فأوراق الفرض
والبديهية لا تتغير.

وعند عقدة MP، سم شجرتي المقدمتين~$T_1,T_2$. جدّد~$T_1$ بعد أن تضم إلى
$H$ جميع ثوابت~$T_2$، ثم جدّد~$T_2$ بعد أن تضم إلى~$H$ جميع ثوابت
الشجرة~$T'_1$ الناتجة. فلا يتساوى ثابت مميز في إحدى الشجرتين بثابت في الأخرى،
ولا يقع في دعمها؛ ويبقى MP صحيحًا، وتصح شروط الجدة بالنسبة إلى الدعم
المجموع.

وعند عقدة~\QR{} ذات الثابت المميز~$a$ والجذر~$!D$، جدّد شجرة المقدمة
أولًا مع تحريم~$H$ و~$a$ وجميع ثوابت~$!D$. ثم اختر
$b\in\mathcal E$ لا يقع في~$H$ ولا في شجرة المقدمة المجددة كلها ولا
في~$!D$، وبدّل الاسمين~$a$ و~$b$ في شجرة المقدمة تبديلًا منتظمًا.
فشرط التطبيق الأول يمنع~$a$ من الدعم ومن~$!B$ و~$!A(x)$، وجدة~$b$
تمنعه من الشجرة؛ فلا يتغير الدعم، ويتحول جذر المقدمة من الحالة ذات~$a$
إلى الحالة ذات~$b$، وتبقى خاتمة عقدة~\QR{} هي~$!D$. وهذا التبديل
تقابل في أسماء الثوابت، فهو ينقل البديهيات إلى أمثلتها، ويحفظ MP والحدود
المغلقة وعقد~\QR{} الداخلة. وقد حظر الاستقراء~$a$ على ثوابتها المميزة،
وحظر اختيار~$b$ على الشجرة كلها، فلا يتحد الثابت المستجد بثابت مميز داخلي.
فنعيد تطبيق~\QR{} بالثابت~$b$؛ والقول واحد في صورتي القاعدة.

فتثبت العبارة الأقوى. خذ فيها~$H$ اتحاد~$F$ وطائفة جميع الثوابت الواقعة
في~$\Sigma_0$، ثم أعد الشجرة إلى متتالية متناهية. فيكون دعمها
$\Sigma\subseteq\Sigma_0$، وتلك~$\Pi'$ المطلوبة.
\end{proof}

\end{document}
